\chapter{绪论：为什么配电网拓扑辨识需要图论}
\label{chap:intro}

\section*{本章目标}
本章建立本书的共同语言：把配电网拓扑辨识表述为图学习问题，区分全节点可观测、仅叶节点可观测与候选图内开关状态辨识三类任务，并说明为什么后续的最小生成树、加性树距离和隐藏树重构不是抽象数学装饰，而是由放射状配电网的物理结构自然引出的工具。后文的树论基础见\cref{chap:trees}，最小生成树理论见\cref{chap:mst}。

\section{拓扑辨识问题的图论表述}
\label{sec:topology_problem}

配电网拓扑辨识的目标，是根据量测数据与先验信息恢复电气连接关系。把母线、配变、分支箱、智能电表等对象抽象为节点，把线路、开关、熔断器或等效馈线段抽象为边，就得到图
\begin{equation}
  G=(V,E),
\end{equation}
其中 $V$ 是节点集合，$E$ 是边集合。若网络在某个运行时段呈放射状，则实际带电拓扑通常是一棵以变电站或台区总表为根的树
\begin{equation}
  T=(V,E_T), \qquad |E_T|=|V|-1 .
\end{equation}

\begin{definition}[配电网拓扑辨识]
给定时序量测 $\mathcal{Y}$、可能的候选边集合 $E_0$ 以及先验信息 $\mathcal{I}$，拓扑辨识是估计实际电气边集
\begin{equation}
  \widehat E_T = \mathcal{A}(\mathcal{Y},E_0,\mathcal{I})
\end{equation}
的过程。若 $\widehat E_T=E_T$，称辨识精确；若二者在隐藏度为 $2$ 节点压缩意义下等价，则称辨识得到等效拓扑。
\end{definition}

这里的“拓扑”不是地理上的相邻，也不是通信链路上的可达，而是电气潮流方程所使用的节点导纳关系。这个区分贯穿全书：\cref{chap:multiview}将专门讨论通信图与电气图不相同的情形。

\section{放射状网、树与有根树}

低压和中压配电网常以放射状方式运行。图论中的树恰好刻画这种结构。

\begin{definition}[树和有根树]
无向连通且无环的图称为树。若指定一个节点 $0\in V$ 为根，则树成为有根树。对非根节点 $i$，从 $0$ 到 $i$ 的唯一路径上与 $i$ 相邻且更靠近根的节点称为 $i$ 的父节点，记为 $\pa(i)$；更远离根的相邻节点称为子节点，集合记为 $\ch(i)$。
\end{definition}

在配电网中，根节点通常对应变电站、馈线首端或台区总表；叶节点通常对应末端用户、电表箱或没有下游支路的节点；内部节点对应分支点、接线箱、开关柜或未装表的中间电气节点。隐藏内部节点是否可恢复，是本书从最小生成树转向隐藏树重构的核心原因，见\cref{chap:hidden_tree}。

\section{末端智能电表量测场景}

实际台区中，智能电表往往部署在用户端。设可观测节点集合为 $L\subseteq V$，通常 $L$ 接近叶节点集合，而内部分支节点 $H=V\setminus L$ 未量测。可得到的时序数据包括有功功率 $P_i(t)$、无功功率 $Q_i(t)$、电压幅值 $V_i(t)$、电流或功率因数等。对仅叶节点可观测场景，算法直接在 $L$ 上建树只会得到一棵连接叶节点的树，而不是原始电气树 $T$。因此，若真实网络含有隐藏分支点，必须估计或重构这些隐藏节点。

\begin{example}[一条馈线的两种观测图]
设真实拓扑为 $0$--$a$--$\{1,2\}$ 与 $a$--$b$--$\{3,4\}$，其中 $0$ 是根，$a,b$ 是分支节点，$1,2,3,4$ 是用户电表。若只有 $1,2,3,4$ 的电压与功率曲线，观测图节点集是 $\{1,2,3,4\}$。即使可以精确得到任意两个电表间的电气距离，也仍需判断 $1,2$ 是否共享隐藏分支点、$3,4$ 是否共享另一个隐藏分支点，这已经超出普通 MST 的输出形式。
\end{example}

\section{多源信息与图关系学习}

拓扑辨识中常见的信息源可以分为四类。

\begin{enumerate}[label=(\arabic*)]
  \item \textbf{电气量测}：$P/Q/V/I$ 时序、停复电事件、电压扰动响应。它们通过潮流方程与电气拓扑相关，是本书\cref{chap:impedance_distance,chap:estimation}的重点。
  \item \textbf{通信量测}：PLC/HPLC 的 RSSI、SNR、RTT、丢包率、路由 hop、载波信道响应等。它们反映通信信道质量，不等同于电气边，见\cref{chap:multiview}。
  \item \textbf{GIS 与台账}：电缆走向、杆塔坐标、箱变归属、开关台账。这些信息给出候选边和物理约束，但可能存在漏录、错录和时效性问题。
  \item \textbf{运行约束}：放射状约束、开关状态约束、潮流可行性和电压上下限约束。它们可用于后处理或联合优化。
\end{enumerate}

因此，一个现代拓扑辨识模型可被看成多视图图关系学习问题：
\begin{equation}
  p_{ij} = \mathbb{P}\{(i,j)\in E_T \mid
  \phi^{V}_{ij},\phi^{PQ}_{ij},\phi^{PLC}_{ij},\phi^{GIS}_{ij}\},
\end{equation}
再把边概率解码为满足放射状约束的图。若概率越大越可信，可用最大生成树；若距离越小越可信，可用最小生成树。

\section{三类拓扑辨识任务}

\subsection{全节点可观测}
若 $V$ 中每个电气节点均有量测，且可构造精确的加性树距离，那么在完全图 $K_V$ 上求 MST 可以恢复真实树。这个结论将在\cref{chap:mst_recovery}严格证明。该情形是 MST 方法最干净、最有理论保证的适用场景。

\subsection{仅叶节点可观测}
若只观测叶节点 $L$，距离矩阵定义在 $L\times L$ 上。真实拓扑含隐藏节点时，$K_L$ 上的 MST 至多产生一棵叶节点树，无法显式恢复隐藏分支点。此时应使用 four-point condition、additive phylogeny、neighbor joining 或 recursive grouping 等隐藏树重构思想，见\cref{chap:four_point,chap:hidden_tree}。

\subsection{候选图内开关状态辨识}
若 GIS 给出候选图 $G_0=(V,E_0)$，任务变为在候选边中选择一组闭合边 $E_T\subseteq E_0$。这类问题可加入放射状约束
\begin{equation}
  |E_T|=|V|-1,\qquad (V,E_T)\ \text{连通},
\end{equation}
也可用混合整数规划、最大生成树、最短路树或贝叶斯图模型求解。

\section{本章小结}

配电网拓扑辨识的数学核心是：用电气、通信和地理信息推断树形或近似树形的电气图。全节点可观测且距离是加性树距离时，MST 有清晰保证；仅叶节点可观测时，隐藏节点导致普通 MST 输出对象不匹配；候选图内开关状态辨识则更像受约束的图选择问题。后续章节将依次补齐树、MST、树距离、阻抗距离与多视图解码的理论基础。

\section*{思考题}
\begin{enumerate}
  \item 为什么“放射状运行”可以抽象为树？若存在联络开关但处于断开状态，它应如何进入图模型？
  \item 若台区总表与所有用户表均可观测，但中间分支箱不可观测，这属于全节点可观测还是叶节点可观测？
  \item 举一个通信图中相邻但电气图中不相邻的例子，并解释可能原因。
  \item 在候选图 $G_0$ 内辨识开关状态时，为什么只根据边概率逐边阈值化可能破坏放射状约束？
\end{enumerate}
